\section{Delay and Mode Selection Theory}
    \label{sec:delay}
    
    \subsection{Circulation Delay as a Physical Primitive}
        
        We introduce the circulation delay as a constitutive element of the theory.
        For a closed vortex filament of effective length $L$ and characteristic propagation speed $v_{\swirlarrow}$, the circulation time is defined as:
        
        \begin{align}
            \tau = \frac{L}{v_{\swirlarrow}}
            \label{eq:delay_time}
        \end{align}
        
        In the special case of a minimal closed loop associated with the vortex core scale, this reduces to:
        
        \begin{align}
            \tau_c = \frac{2\pi r_c}{v_{\swirlarrow}} = \frac{2\pi}{\omega_c}
            \label{eq:compton_delay}
        \end{align}
        
        This establishes a direct identification between the circulation delay and the inverse Compton frequency.
        
        \textbf{[DERIVED]}: The intrinsic delay of the vortex loop is identical to the Compton period.
    
    \subsection{Delayed Self-Interaction Dynamics}
        
        We model the phase $\phi(t)$ of a circulating excitation as a delayed self-interacting oscillator:
        
        \begin{align}
            \dot{\phi}(t) = \omega_0 + \kappa \sin\big(\phi(t - \tau) - \phi(t)\big)
            \label{eq:delay_phase}
        \end{align}
        
        This equation represents a minimal effective description of a circulating field with finite propagation time.
        
        \textbf{[ORTHODOX]}: Delay-differential equations of this type are well-established in feedback systems \cite{SST_delay_basic}.
    
    \subsection{Emergence of Discrete Stable Modes}
        
        We seek steady-state solutions of the form:
        
        \begin{align}
            \phi(t) = \Omega t
        \end{align}
        
        Substituting into Eq.~\ref{eq:delay_phase} yields:
        
        \begin{align}
            \Omega = \omega_0 + \kappa \sin(\Omega \tau)
            \label{eq:mode_condition}
        \end{align}
        
        This transcendental equation admits a discrete set of stable solutions $\Omega_n$ determined by the delay parameter $\tau$.
        
        \textbf{[DERIVED]}: Discrete mode selection arises from temporal nonlocality alone.
    
    \subsection{Delay-Induced Quantization Mechanism}
        
        Unlike conventional quantum mechanics, where discreteness is postulated, here it emerges dynamically:
        
        \begin{itemize}
            \item No boundary condition required
            \item No operator quantization assumed
            \item Discreteness arises from stability selection
        \end{itemize}
        
        \textbf{[SPECULATIVE]}: Particle spectra may be interpreted as delay-selected stable circulation modes.
    
    \subsection{Extension to Nonlinear Schrödinger Dynamics}
        
        To connect with field-theoretic descriptions, we extend the delay mechanism to a nonlinear Schrödinger framework:
        
        \begin{align}
            i \hbar \partial_t \psi = -\frac{\hbar^2}{2m} \nabla^2 \psi + F_{\text{delay}}[\psi(t - \tau)]
        \end{align}
        
        where the delayed functional introduces temporal nonlocality into the field evolution.
        
        \textbf{[SPECULATIVE]}: This provides a bridge between classical delay systems and quantum wave dynamics.
    
    \subsection{Physical Interpretation}
        
        The delay $\tau$ acts as:
        
        \begin{itemize}
            \item a memory kernel
            \item a nonlocal temporal coupling
            \item a selection mechanism for stable modes
        \end{itemize}
        
        Thus:
        
        \begin{center}
            \textit{Delay $\Rightarrow$ Stability $\Rightarrow$ Discreteness}
        \end{center}
        
        This mechanism forms one of the central pillars of the SST canon.